%-------------------------
% Resume in Latex
% Author : Jake Gutierrez
% Based off of: https://github.com/sb2nov/resume
% License : MIT
%------------------------

\documentclass[letterpaper,11pt]{article}

\usepackage{latexsym}
\usepackage[empty]{fullpage}
\usepackage{titlesec}
\usepackage{marvosym}
\usepackage[usenames,dvipsnames]{color}
\usepackage{verbatim}
\usepackage{enumitem}
\usepackage[hidelinks]{hyperref}
\usepackage{fancyhdr}
\usepackage[english]{babel}
\usepackage{tabularx}
\input{glyphtounicode}


%----------FONT OPTIONS----------
% sans-serif
% \usepackage[sfdefault]{FiraSans}
% \usepackage[sfdefault]{roboto}
% \usepackage[sfdefault]{noto-sans}
% \usepackage[default]{sourcesanspro}

% serif
% \usepackage{CormorantGaramond}
% \usepackage{charter}


\pagestyle{fancy}
\fancyhf{} % clear all header and footer fields
\fancyfoot{}
\renewcommand{\headrulewidth}{0pt}
\renewcommand{\footrulewidth}{0pt}

% Adjust margins
\addtolength{\oddsidemargin}{-0.5in}
\addtolength{\evensidemargin}{-0.5in}
\addtolength{\textwidth}{1in}
\addtolength{\topmargin}{-.5in}
\addtolength{\textheight}{1.0in}

\urlstyle{same}

\raggedbottom
\raggedright
\setlength{\tabcolsep}{0in}

% Sections formatting
\titleformat{\section}{
  \vspace{-4pt}\scshape\raggedright\large
}{}{0em}{}[\color{black}\titlerule \vspace{-5pt}]

% Ensure that generate pdf is machine readable/ATS parsable
\pdfgentounicode=1

%-------------------------
% Custom commands
\newcommand{\resumeItem}[1]{
  \item\small{
    {#1 \vspace{-2pt}}
  }
}

\newcommand{\resumeSubheading}[4]{
  \vspace{-2pt}\item
    \begin{tabular*}{0.97\textwidth}[t]{l@{\extracolsep{\fill}}r}
      \textbf{#1} & #2 \\
      \textit{\small#3} & \textit{\small #4} \\
    \end{tabular*}\vspace{-7pt}
}

\newcommand{\resumeSubSubheading}[2]{
    \item
    \begin{tabular*}{0.97\textwidth}{l@{\extracolsep{\fill}}r}
      \textit{\small#1} & \textit{\small #2} \\
    \end{tabular*}\vspace{-7pt}
}

\newcommand{\resumeProjectHeading}[2]{
    \item
    \begin{tabular*}{0.97\textwidth}{l@{\extracolsep{\fill}}r}
      \small#1 & #2 \\
    \end{tabular*}\vspace{-7pt}
}

\newcommand{\resumeSubItem}[1]{\resumeItem{#1}\vspace{-4pt}}

\renewcommand\labelitemii{$\vcenter{\hbox{\tiny$\bullet$}}$}

\newcommand{\resumeSubHeadingListStart}{\begin{itemize}[leftmargin=0.15in, label={}]}
\newcommand{\resumeSubHeadingListEnd}{\end{itemize}}
\newcommand{\resumeItemListStart}{\begin{itemize}}
\newcommand{\resumeItemListEnd}{\end{itemize}\vspace{-5pt}}

%-------------------------------------------
%%%%%%  RESUME STARTS HERE  %%%%%%%%%%%%%%%%%%%%%%%%%%%%


\begin{document}

%----------HEADING----------
% \begin{tabular*}{\textwidth}{l@{\extracolsep{\fill}}r}
%   \textbf{\href{http://sourabhbajaj.com/}{\Large Sourabh Bajaj}} & Email : \href{mailto:sourabh@sourabhbajaj.com}{sourabh@sourabhbajaj.com}\\
%   \href{http://sourabhbajaj.com/}{http://www.sourabhbajaj.com} & Mobile : +1-123-456-7890 \\
% \end{tabular*}

\begin{center}
    \textbf{\Huge \scshape Yisheng Gong} \\ \vspace{1pt}
    \small +1 314-904-6457(US) $|$ \href{mailto:yishenggong@9437.com}{{yishenggong9437@gmail.com}} $|$ \href{https://www.linkedin.com/in/yisheng-gong1997}{{linkedin}} $|$
    \href{https://www.github.com/gongyisheng}{{github}} $|$ \href{https://blog.yellowday.day}{{website}}
\end{center}


%-----------EDUCATION-----------
\section{Education}
  \resumeSubHeadingListStart
    \resumeSubheading
      {Washington University in St.Louis}{St.Louis, MO}
      {Master of Science in Business Analytics}{Jan 2021 -- May 2022}
    \resumeSubheading
      {Fudan University}{Shanghai, China}
      {Bachelor of Science in Chemistry }{Sept 2015 -- June 2019}
  \resumeSubHeadingListEnd

  %
%-----------PROGRAMMING SKILLS-----------
\section{Technical Skills}
 \begin{itemize}[leftmargin=0.15in, label={}]
    \small{\item{
     \textbf{Languages}{: Python, Go, Bash} \\
     \textbf{Infra}{: Kubernetes, Docker, Prometheus, Grafana, CI/CD} \\
     \textbf{AI/LLM}{: LLM inference optimization, KV Cache, SFT, RL, model serving} \\
     \textbf{ML Frameworks}{: PyTorch, vLLM, SGLang, Huggingface Transformers, veRL (contributor), miles (contributor)} \\
     \textbf{Cloud}{: AWS (EC2/S3)} \\
    }}
 \end{itemize}


%-----------EXPERIENCE-----------
\section{Professional Experience}
  \resumeSubHeadingListStart
    \resumeSubheading
      {Open Source Contributor}{Dec 2025 -- }
      {Miles / SGLang}{}
      \resumeItemListStart
        \resumeItem{Participated in LoRA support for the Megatron backend in the Miles framework, including MoE layer LoRA support for GPT-OSS model. Implemented LoRA weights synchronization to the SGLang inference engine and colocate model feature with train/rollout offloading, enabling efficient adapter serving and GPU resource sharing in distributed training-inference pipelines.}
        \resumeItem{Fixed a kernel bug in SGLang that caused garbage output in the batch invariant operation kernel on Ada/Ampere GPUs.}
      \resumeItemListEnd
    \resumeSubheading
      {Software Engineer, Machine Learning}{June 2022 -- }
      {Yipitdata, Edison Platform Team}{Mountain View, CA, US}
      \resumeItemListStart
      \resumeItem{Designed and implemented an automated AI tagging agent for vendor matching. Optimized agent performance with context engineering, including dynamic context onload/offload, majority-k voting for stability, and model routing for cost reduction. Achieved 95\% accuracy at \$5k per 1M data records, slightly surpassing human taggers while saving 7 headcounts.}
      \resumeItem{Developed a multi-turn SFT + search-R1 like Agentic RL post-training pipeline for a 14B LLM. Use GPT-5-mini reasoning traces for SFT and a repeated rollout sampling strategy to remove overly hard or low-variance cases that weaken reward gradients. Model trained on 4xL40S instance. Delivered a 30\% (60\% to 90\%) accuracy jump on identifying ambiguous case, got same performance as GPT-5 class baselines.}
      \resumeItem{Optimized core service infrastructure using profiling and network packet capture to identify bottlenecks. Achieved 55\% higher throughput, 70\% lower cost, and latency reductions. Designed redis client cache component (-90\% I/O) and event-driven MySQL binlog to Redis Streams sync system for real-time data delivery.}
      \resumeItem{Worked with DevOps team to design and implement an HPA (Horizontal Pod Autoscaler) system based on KEDA, enabling auto-scaling of services based on custom metrics.}
      \resumeItemListEnd
    \resumeSubheading
      {Software Engineer, Part Time}{Jan 2022 -- May 2022}
      {Massachusetts Institute of Technology, ChemE, CRIPT Team}{Remote, US}
    \resumeSubheading
      {Data Engineer Intern}{Jan 2021 -- May 2021}
      {Bytedance(TikTok), Group of Douyin E-commerce, Commercialization Division}{Beijing, China}
  \resumeSubHeadingListEnd
%-------------------------------------------
\end{document}

